\chapter{Program Architecture for a Full-Day Internship}

\section{Why This Is a 10-Week Full-Day Program}
The program is not 10 weeks because students need 10 weeks to learn a few R functions. It is 10 weeks because the expected outcome is a research product with defensible evidence, reusable code, figures, interpretation, and a handoff package. A full-day internship gives students enough time to repeat the cycle that real research requires: learn, try, fail, diagnose, revise, document, and communicate.

\begin{programbox}{Full-Day Assumption}
Each week assumes roughly 30 to 35 focused hours. A typical day includes a short technical lesson, guided work, independent project time, dsDNA Core or peer feedback, and written documentation. Students who only attend part-time can still use the program; in that format the number of exercises is reduced rather than reducing the quality standard.
\end{programbox}

\begin{goldbox}{Weekend Asynchronous Expectation}
The program uses all 70 calendar days. Weekend work is lighter than weekday work, but it is required. Weekend tasks are designed for one to three focused hours and usually include reading annotation, rerunning a small script, updating the research log, revising a worksheet, reviewing a peer artifact, or preparing the next week's files. Weekend work should never require live database access unless that activity has been explicitly planned and cached examples are available as a fallback.
\end{goldbox}

\section{Daily Work Blocks}
\begin{tabularx}{\textwidth}{p{0.22\textwidth}Y}
\toprule
\textbf{Block} & \textbf{Purpose}\\
\midrule
Morning concept block & Introduce the scientific idea, database concept, coding pattern, or interpretation principle needed that day.\\
Guided technical lab & Run code, inspect outputs, troubleshoot errors, and learn the expected workflow with a small example.\\
Independent project block & Apply the day's method to the student's own project or assigned case study.\\
Evidence and writing block & Convert analysis into a research log entry, methods text, figure caption, or decision memo.\\
dsDNA Core checkpoint & Review a concrete artifact and decide whether the student can move forward or must revise.\\
\bottomrule
\end{tabularx}

\section{Weekday and Weekend Time Design}
\begin{longtable}{p{0.24\textwidth}p{0.24\textwidth}p{0.42\textwidth}}
\toprule
\textbf{Day type} & \textbf{Typical time} & \textbf{Expected work}\\
\midrule
\endfirsthead
\toprule
\textbf{Day type} & \textbf{Typical time} & \textbf{Expected work}\\
\midrule
\endhead
Weekday concept and guided work & 2 to 3 hours & Reading discussion, demonstration, code walk-through, or source interpretation exercise.\\
Weekday independent project work & 2 to 3 hours & Student applies the day's method to their own project or a case dataset and saves an artifact.\\
Weekday evidence and writing & 1 to 1.5 hours & Research log, method text, figure caption, worksheet, claim memo, or review response.\\
Weekday review checkpoint & 20 to 45 minutes & Peer or dsDNA Core review of a specific artifact and next-step decision.\\
Weekend asynchronous work & 1 to 3 hours & Reading annotation, rerun, project cleanup, source audit, reflection, peer review, or revision.\\
\bottomrule
\end{longtable}

\section{Product Ladder}
The internship is organized as a product ladder. Each week creates an artifact that is needed later.

\begin{longtable}{p{0.12\textwidth}p{0.25\textwidth}p{0.53\textwidth}}
\toprule
\textbf{Week} & \textbf{Artifact} & \textbf{Why it matters}\\
\midrule
\endfirsthead
\toprule
\textbf{Week} & \textbf{Artifact} & \textbf{Why it matters}\\
\midrule
\endhead
1 & Project brief & Prevents a broad interest from becoming an untestable analysis.\\
2 & Reproducible project & Gives the student a research workspace that can be audited and rerun.\\
3 & Database evidence map & Separates identity, source coverage, biological context, and uncertainty.\\
4 & Core GC-MS workflow & Establishes how tentative chemical evidence enters the project.\\
5 & Enrichment result & Creates the main source-backed chemical evidence object.\\
6 & Grouping dictionary & Converts source-backed fields into analysis-ready variables.\\
7 & Figure set & Reveals patterns and missing data before interpretation hardens.\\
8 & Evidence memo & Connects database evidence to peer-reviewed context.\\
9 & Draft report & Forces the methods, results, figures, and limitations into a coherent argument.\\
10 & Final package & Leaves behind work that the dsDNA Core team, a collaborator, or a future student can use.\\
\bottomrule
\end{longtable}

\section{Skill Progression}
Students move through four levels:

\begin{enumerate}
  \item \textbf{Operator}: can run a provided script and identify the output.
  \item \textbf{Inspector}: can check whether an output is complete, duplicated, missing, or suspicious.
  \item \textbf{Analyst}: can transform validated outputs into variables, figures, and tables.
  \item \textbf{Researcher}: can defend interpretation with source evidence, literature, uncertainty, and next steps.
\end{enumerate}

\section{Research Integrity Rules}
\begin{checklistbox}{Rules Students Must Follow}
\begin{itemize}
  \item Never present a tentative chemical match as confirmed identity unless standards or equivalent confirmatory evidence are available.
  \item Never hide unmatched chemicals or missing source coverage.
  \item Never copy a database sentence into a result and treat it as an analyzed variable.
  \item Never store tokens, passwords, private data, or unpublished collaborator data in a repository.
  \item Never change raw data manually; all cleaning must be scripted or documented.
  \item Never make a figure that cannot be regenerated from code.
\end{itemize}
\end{checklistbox}

\section{Program Decision Gates}
Decision gates keep the program rigorous:

\begin{tabularx}{\textwidth}{p{0.22\textwidth}Y}
\toprule
\textbf{Gate} & \textbf{Student must show}\\
\midrule
Scope gate & A feasible question, reviewed chemical list, and documented risk note.\\
Reproducibility gate & A project folder with relative paths, scripts, logs, and no secrets.\\
Evidence gate & Source coverage, validation results, and evidence rows for important traits.\\
Variable gate & A grouping dictionary with source table, field, rule, allowed values, and caveat.\\
Interpretation gate & A claim supported by a figure, table, database field, and literature source.\\
Handoff gate & Final package can be understood by someone who was not present during the internship.\\
\bottomrule
\end{tabularx}

\section{Expected Weekly Time Budget}
\begin{tabularx}{\textwidth}{p{0.35\textwidth}p{0.18\textwidth}Y}
\toprule
\textbf{Activity} & \textbf{Hours/week} & \textbf{Notes}\\
\midrule
Technical instruction and guided lab & 6--8 & Short lessons followed by hands-on work.\\
Independent project work & 10--12 & Student applies the workflow to their question or case study.\\
Reading and annotation & 3--5 & Database docs, methods papers, and project-specific literature.\\
Writing and documentation & 4--6 & Research log, methods, captions, interpretation, and handoff notes.\\
dsDNA Core and peer review & 3--4 & Weekly review, code walk-throughs, and figure critique.\\
Troubleshooting and revision & 4--6 & Expected time for errors, reruns, cleanup, and rethinking.\\
\bottomrule
\end{tabularx}

\section{What Students Should Be Told Up Front}
The internship is not a race to get a colorful heatmap. A student who finds that half their chemicals have weak source coverage has not failed. That is a useful result if it is documented clearly. A student who discovers that a project question is too broad has not failed. That is research judgment. The standard is not that every project produces a dramatic biological conclusion. The standard is that every project produces a trustworthy research package.

\begin{checklistbox}{Student Briefing Before Week 1}
\begin{itemize}
  \item The program is a full-day research experience with required weekend asynchronous work.
  \item The final product is a reproducible research package, not only a presentation.
  \item uafR helps organize and enrich chemical evidence, but it does not confirm identities by itself.
  \item Public database evidence must be interpreted with source coverage, provenance, and caveats.
  \item Research logs are evaluated because they preserve decisions and make the project reusable.
  \item Students may bring their own interests, but each project must pass the same evidence and reproducibility gates.
\end{itemize}
\end{checklistbox}

\section{What the Program Should Prepare Before Week 1}
\begin{checklistbox}{Preparation Checklist}
\begin{itemize}
  \item Current private uafR student bundle or verified installation route.
  \item Example chemical list and case dataset for students without immediate project data.
  \item Printed or digital worksheets for project brief, chemical intake, database evidence map, and research log.
  \item A clear data classification note for public, private, unpublished, and restricted materials.
  \item A shared location for program-safe examples, not for private raw collaborator data unless access is approved.
  \item A review calendar with Week 1 scope gate, Week 2 reproducibility gate, Week 5 validation gate, Week 6 variable gate, Week 8 interpretation gate, and Week 10 handoff gate.
  \item A backup plan for days when live public database services are unavailable.
\end{itemize}
\end{checklistbox}

\section{How Work Accumulates}
Each week reuses the previous week's artifact. The project brief controls the chemical intake table. The chemical intake table controls the enrichment query. The saved enrichment result controls source coverage and trait extraction. The grouping dictionary controls the matrix. The matrix controls figures. Figures control results paragraphs. Results paragraphs control the discussion and presentation. This chain prevents the final report from becoming a collection of disconnected outputs.
